\begin{frame}{이 장의 한 줄: 정확도의 대가, 그리고 되찾기}
  \begin{block}{한 줄}
    결정 트리는 ``사람이 읽을 수 있는 if-then 규칙의 나무''이면서 동시에,
    ``얼마나 잘 나누는가''를 \textbf{정확히 수식으로 잴 수 있는}
    알고리즘이다.
  \end{block}
  \vskip0.8em
  \begin{itemize}
    \item 트리를 \textbf{여러 개 합치면}(랜덤 포레스트, GBDT)
      \textbf{정확도는 올라가지만} 그 읽기 쉬움을 \textbf{잃는다}
    \item \kb{SHAP}은 그 대가로 잃은 설명가능성을,
      \textbf{게임이론이라는 전혀 다른 도구}로 \textbf{되찾아오는} 시도
  \end{itemize}
  \vskip0.8em
  \begin{itemize}
    \item \textbf{7.1} — 좋은 질문 = \textbf{지니불순도 감소량(정보이득)
      최대}. U자 곡선: \textbf{test 정점}이 최적 깊이.
    \item \textbf{7.2} — \textbf{독립적 다수결}. 두 무작위성이 $\rho$를
      줄여 \textbf{분산} 감소. 계단 $\to$ 부드러운 곡선.
    \item \textbf{7.3} — \textbf{순차적 잔차 추격} = 함수 공간의
      경사하강. 더 정밀한 대신 \textbf{early stopping}이 필수.
      SHAP(가산성)으로 예측 하나를 \textbf{정확히} 분해.
  \end{itemize}
  \vskip0.8em
  {\footnotesize 7.1절의 ``트리 하나를 잘 나누는 원리''는 이 장의
  \textbf{전부에서 재사용}되었다.}
\end{frame}
